\chapter{Hemorrhoidectomy}

\section*{What Is This Test or Procedure?}
Hemorrhoidectomy is surgical removal of enlarged hemorrhoidal tissue, usually for severe or recurrent hemorrhoids.

\section*{Why This Test or Procedure Is Ordered}
It may be considered for large prolapsing hemorrhoids, persistent bleeding or pain, strangulation, or symptoms that do not improve with dietary measures and office treatments.

\section*{How This Test or Procedure Is Performed}
Under anesthesia, the surgeon excises selected hemorrhoidal columns and controls the blood supply. Open, closed, or stapled techniques may be used depending on the anatomy.

\section*{Risks and Recovery}
Pain, bleeding, urinary retention, infection, anal stenosis, recurrence, and temporary bowel-control problems are possible. Stool-softening measures, fluids, local care, and activity guidance are important during recovery.

\section*{References}
\begin{enumerate}
  \item MedlinePlus. Hemorrhoidectomy. U.S. National Library of Medicine; 2025.
  \item Current Medical Diagnosis and Treatment. McGraw-Hill; 2025.
\end{enumerate}
\clearpage
